% Generated draft: requires mathematical and source review.
\documentclass[11pt]{article}
\usepackage[margin=25mm]{geometry}
\usepackage{fontspec}
\setmainfont{Latin Modern Roman}
\usepackage{amsmath,amssymb,mathtools,unicode-math}
\setmathfont{Latin Modern Math}
\usepackage{enumitem,needspace,xurl,hyperref}
\setlength{\parindent}{0pt}
\setlength{\parskip}{5pt}
\setlength{\emergencystretch}{3em}
\newcommand{\N}{\mathbb{N}}
\newcommand{\Z}{\mathbb{Z}}
\newcommand{\Q}{\mathbb{Q}}
\newcommand{\R}{\mathbb{R}}
\newcommand{\C}{\mathbb{C}}
\newcommand{\F}{\mathbb{F}}
\newcommand{\A}{\mathbb{A}}
\newcommand{\PP}{\mathbb{P}}
\newcommand{\E}{\mathbb{E}}
\newcommand{\eps}{\varepsilon}
\newcommand{\dd}{\,\mathrm{d}}
\begin{document}
\section*{UnsolvedMath ID: 800846}
\textbf{Source problem code:} catalogue entry
\subsection{Definitions and mathematical statement}
For each integer \(k\geq 0\), define the Sylvester Hadamard matrix recursively by
\[
H_0=(1),\qquad H_{k+1}=\begin{pmatrix}H_k&H_k\\H_k&-H_k\end{pmatrix}.
\]
Thus \(H_k\in\{-1,1\}^{2^k\times 2^k}\). For a real \(m\times n\) matrix \(A=(a_{ij})\), interpret its discrepancy as the unnormalized sign discrepancy
\[
\operatorname{disc}(A):=\min_{x\in\{-1,1\}^n}\|Ax\|_\infty
=\min_{x\in\{-1,1\}^n}\max_{1\leq i\leq m}\left|\sum_{j=1}^n a_{ij}x_j\right|.
\]
The supplied excerpt does not define \(\operatorname{disc}\) and refers to preceding material for the Sylvester construction. Consequently, the following is a conditional formalization using the explicitly stated conventions; their agreement with the full source requires review.

The supplied Conjecture 12 asserts that, for every odd integer \(k\geq 1\),
\[
\operatorname{disc}(H_k)=\sqrt{2}\sqrt{2^k}.
\]
Equivalently, for each such \(k\), there exists \(x\in\{-1,1\}^{2^k}\) such that every coordinate of \(H_kx\) has absolute value at most \(\sqrt{2}\sqrt{2^k}\), and every \(x\in\{-1,1\}^{2^k}\) has at least one coordinate of \(H_kx\) with absolute value at least \(\sqrt{2}\sqrt{2^k}\).
\subsection{Short English statement}
Under the stated discrepancy convention, the conjecture says that for every positive odd \(k\), the smallest possible largest absolute row sum of the Sylvester matrix after assigning signs to its columns is exactly \(\sqrt{2}\sqrt{2^k}\). The supplied registry labels the problem open; that status and the discrepancy convention have not been independently verified.
\subsection{Sources}
\begin{itemize}
\item Supplied local UnsolvedMath registry record, permanent ID \(800846\), source problem code \texttt{catalogue entry}, titled ``Randomstrasse101 Open Problems (2024) --- Conjecture 12.'' The record attributes the problem to Bandeira et al., gives proposed year 2025, and labels its status as open.
\item Bandeira et al., \emph{Randomstrasse101: Open Problems of 2024} (2025), Conjecture 12. Supplied source URL: \url{https://arxiv.org/abs/2504.20539}. The record reports access on 2026-10-16 and extraction by \texttt{arxiv\_latex}; these are supplied provenance details, not an independent verification or a claim of access here.
\item Supplied upstream revision: \texttt{d371ee917ac9647fbb06cd6fc529deec6f3dda82}.
\end{itemize}
\end{document}
